\documentclass[12pt,a4paper]{exam}
\usepackage[utf8]{inputenc}
\usepackage{amsmath,amssymb,amsthm}
\usepackage[margin=1in]{geometry}
\usepackage{graphicx}
\usepackage[colorlinks=true]{hyperref}
\pagestyle{headandfoot}
\firstpageheader{MIT OpenCourseWare}{}{\textbf{EXTREMELY HARD -- MIT LEVEL}}
\runningheader{MIT OpenCourseWare}{}{\textbf{EXTREMELY HARD -- MIT LEVEL}}
\firstpagefooter{}{Page \thepage}{}
\runningfooter{}{Page \thepage}{}

\begin{document}

\begin{center}
{\LARGE \textbf{MIT OpenCourseWare}}\\5.111\\Principles of Chemical Science
\vspace{0.5cm}
{5.111} --- {Principles of Chemical Science}
\vspace{0.3cm}
May 2026
\vspace{0.3cm}
\textbf{EXTREMELY HARD -- MIT LEVEL}
\end{center}

\begin{center}
\textbf{Time Limit: 3 hours}
\end{center}

\begin{center}
\textbf{Total Marks: 100}
\end{center}

\vspace{0.5cm}

\begin{questions}

\question[2]
Analyze the limitations of standard approaches to Atomic Theory when applied to edge cases in 5.111. Construct explicit counterexamples showing where intuitive reasoning fails.

\vspace{0.3cm}

\question[5]
Construct an original proof-level problem involving Chemical Bonding for 5.111 (Principles of Chemical Science) and provide a complete solution. Your problem should be at the level of a graduate qualifying exam.

\vspace{0.3cm}

\question[4]
Analyze the limitations of standard approaches to Thermodynamics when applied to edge cases in 5.111. Construct explicit counterexamples showing where intuitive reasoning fails.

\vspace{0.3cm}

\question[6]
Construct an original proof-level problem involving Kinetics for 5.111 (Principles of Chemical Science) and provide a complete solution. Your problem should be at the level of a graduate qualifying exam.

\vspace{0.3cm}

\question[2]
Construct an original proof-level problem involving Equilibrium for 5.111 (Principles of Chemical Science) and provide a complete solution. Your problem should be at the level of a graduate qualifying exam.

\vspace{0.3cm}

\question[5]
Analyze the limitations of standard approaches to Acid-Base Chemistry when applied to edge cases in 5.111. Construct explicit counterexamples showing where intuitive reasoning fails.

\vspace{0.3cm}

\question[7]
Prove the fundamental theorem concerning Atomic Theory in the context of 5.111 (Principles of Chemical Science). State the theorem precisely, provide a complete proof, and discuss three important corollaries.

\vspace{0.3cm}

\question[3]
Prove the fundamental theorem concerning Chemical Bonding in the context of 5.111 (Principles of Chemical Science). State the theorem precisely, provide a complete proof, and discuss three important corollaries.

\vspace{0.3cm}

\question[2]
Analyze the limitations of standard approaches to Thermodynamics when applied to edge cases in 5.111. Construct explicit counterexamples showing where intuitive reasoning fails.

\vspace{0.3cm}

\question[3]
Construct an original proof-level problem involving Kinetics for 5.111 (Principles of Chemical Science) and provide a complete solution. Your problem should be at the level of a graduate qualifying exam.

\vspace{0.3cm}

\question[5]
Prove the fundamental theorem concerning Equilibrium in the context of 5.111 (Principles of Chemical Science). State the theorem precisely, provide a complete proof, and discuss three important corollaries.

\vspace{0.3cm}

\question[5]
Analyze the limitations of standard approaches to Acid-Base Chemistry when applied to edge cases in 5.111. Construct explicit counterexamples showing where intuitive reasoning fails.

\vspace{0.3cm}

\question[6]
Analyze the limitations of standard approaches to Atomic Theory when applied to edge cases in 5.111. Construct explicit counterexamples showing where intuitive reasoning fails.

\vspace{0.3cm}

\question[5]
Prove the fundamental theorem concerning Chemical Bonding in the context of 5.111 (Principles of Chemical Science). State the theorem precisely, provide a complete proof, and discuss three important corollaries.

\vspace{0.3cm}

\question[4]
Prove the fundamental theorem concerning Thermodynamics in the context of 5.111 (Principles of Chemical Science). State the theorem precisely, provide a complete proof, and discuss three important corollaries.

\vspace{0.3cm}

\question[2]
Analyze the limitations of standard approaches to Kinetics when applied to edge cases in 5.111. Construct explicit counterexamples showing where intuitive reasoning fails.

\vspace{0.3cm}

\question[3]
Construct an original proof-level problem involving Equilibrium for 5.111 (Principles of Chemical Science) and provide a complete solution. Your problem should be at the level of a graduate qualifying exam.

\vspace{0.3cm}

\question[2]
Prove the fundamental theorem concerning Acid-Base Chemistry in the context of 5.111 (Principles of Chemical Science). State the theorem precisely, provide a complete proof, and discuss three important corollaries.

\vspace{0.3cm}

\question[3]
Analyze the limitations of standard approaches to Atomic Theory when applied to edge cases in 5.111. Construct explicit counterexamples showing where intuitive reasoning fails.

\vspace{0.3cm}

\question[6]
Prove the fundamental theorem concerning Chemical Bonding in the context of 5.111 (Principles of Chemical Science). State the theorem precisely, provide a complete proof, and discuss three important corollaries.

\vspace{0.3cm}

\question[5]
Construct an original proof-level problem involving Thermodynamics for 5.111 (Principles of Chemical Science) and provide a complete solution. Your problem should be at the level of a graduate qualifying exam.

\vspace{0.3cm}

\question[5]
Prove the fundamental theorem concerning Kinetics in the context of 5.111 (Principles of Chemical Science). State the theorem precisely, provide a complete proof, and discuss three important corollaries.

\vspace{0.3cm}

\question[2]
Prove the fundamental theorem concerning Equilibrium in the context of 5.111 (Principles of Chemical Science). State the theorem precisely, provide a complete proof, and discuss three important corollaries.

\vspace{0.3cm}

\question[5]
Construct an original proof-level problem involving Acid-Base Chemistry for 5.111 (Principles of Chemical Science) and provide a complete solution. Your problem should be at the level of a graduate qualifying exam.

\vspace{0.3cm}

\question[3]
Analyze the limitations of standard approaches to Atomic Theory when applied to edge cases in 5.111. Construct explicit counterexamples showing where intuitive reasoning fails.

\vspace{0.3cm}

\end{questions}

\newpage
\section*{Solutions}

\textbf{Solution 1:}

The edge case analysis reveals the subtle assumptions underlying standard treatments of Atomic Theory. By constructing explicit counterexamples, we see that the usual hypotheses cannot be weakened without loss of the theorem's conclusion.

\vspace{0.5cm}

\textbf{Solution 2:}

Solution approach: First recall the definitions and key theorems related to Chemical Bonding from 5.111. Then apply the appropriate techniques to construct a rigorous argument. The solution must be self-contained and reference only the material from the course.

\vspace{0.5cm}

\textbf{Solution 3:}

The edge case analysis reveals the subtle assumptions underlying standard treatments of Thermodynamics. By constructing explicit counterexamples, we see that the usual hypotheses cannot be weakened without loss of the theorem's conclusion.

\vspace{0.5cm}

\textbf{Solution 4:}

Solution approach: First recall the definitions and key theorems related to Kinetics from 5.111. Then apply the appropriate techniques to construct a rigorous argument. The solution must be self-contained and reference only the material from the course.

\vspace{0.5cm}

\textbf{Solution 5:}

Solution approach: First recall the definitions and key theorems related to Equilibrium from 5.111. Then apply the appropriate techniques to construct a rigorous argument. The solution must be self-contained and reference only the material from the course.

\vspace{0.5cm}

\textbf{Solution 6:}

The edge case analysis reveals the subtle assumptions underlying standard treatments of Acid-Base Chemistry. By constructing explicit counterexamples, we see that the usual hypotheses cannot be weakened without loss of the theorem's conclusion.

\vspace{0.5cm}

\textbf{Solution 7:}

The proof follows from the definitions and properties established in 5.111. The key insight is to apply the relevant theorem and verify the hypotheses hold. Detailed solution depends on the specific formulation of the theorem.

\vspace{0.5cm}

\textbf{Solution 8:}

The proof follows from the definitions and properties established in 5.111. The key insight is to apply the relevant theorem and verify the hypotheses hold. Detailed solution depends on the specific formulation of the theorem.

\vspace{0.5cm}

\textbf{Solution 9:}

The edge case analysis reveals the subtle assumptions underlying standard treatments of Thermodynamics. By constructing explicit counterexamples, we see that the usual hypotheses cannot be weakened without loss of the theorem's conclusion.

\vspace{0.5cm}

\textbf{Solution 10:}

Solution approach: First recall the definitions and key theorems related to Kinetics from 5.111. Then apply the appropriate techniques to construct a rigorous argument. The solution must be self-contained and reference only the material from the course.

\vspace{0.5cm}

\textbf{Solution 11:}

The proof follows from the definitions and properties established in 5.111. The key insight is to apply the relevant theorem and verify the hypotheses hold. Detailed solution depends on the specific formulation of the theorem.

\vspace{0.5cm}

\textbf{Solution 12:}

The edge case analysis reveals the subtle assumptions underlying standard treatments of Acid-Base Chemistry. By constructing explicit counterexamples, we see that the usual hypotheses cannot be weakened without loss of the theorem's conclusion.

\vspace{0.5cm}

\textbf{Solution 13:}

The edge case analysis reveals the subtle assumptions underlying standard treatments of Atomic Theory. By constructing explicit counterexamples, we see that the usual hypotheses cannot be weakened without loss of the theorem's conclusion.

\vspace{0.5cm}

\textbf{Solution 14:}

The proof follows from the definitions and properties established in 5.111. The key insight is to apply the relevant theorem and verify the hypotheses hold. Detailed solution depends on the specific formulation of the theorem.

\vspace{0.5cm}

\textbf{Solution 15:}

The proof follows from the definitions and properties established in 5.111. The key insight is to apply the relevant theorem and verify the hypotheses hold. Detailed solution depends on the specific formulation of the theorem.

\vspace{0.5cm}

\textbf{Solution 16:}

The edge case analysis reveals the subtle assumptions underlying standard treatments of Kinetics. By constructing explicit counterexamples, we see that the usual hypotheses cannot be weakened without loss of the theorem's conclusion.

\vspace{0.5cm}

\textbf{Solution 17:}

Solution approach: First recall the definitions and key theorems related to Equilibrium from 5.111. Then apply the appropriate techniques to construct a rigorous argument. The solution must be self-contained and reference only the material from the course.

\vspace{0.5cm}

\textbf{Solution 18:}

The proof follows from the definitions and properties established in 5.111. The key insight is to apply the relevant theorem and verify the hypotheses hold. Detailed solution depends on the specific formulation of the theorem.

\vspace{0.5cm}

\textbf{Solution 19:}

The edge case analysis reveals the subtle assumptions underlying standard treatments of Atomic Theory. By constructing explicit counterexamples, we see that the usual hypotheses cannot be weakened without loss of the theorem's conclusion.

\vspace{0.5cm}

\textbf{Solution 20:}

The proof follows from the definitions and properties established in 5.111. The key insight is to apply the relevant theorem and verify the hypotheses hold. Detailed solution depends on the specific formulation of the theorem.

\vspace{0.5cm}

\textbf{Solution 21:}

Solution approach: First recall the definitions and key theorems related to Thermodynamics from 5.111. Then apply the appropriate techniques to construct a rigorous argument. The solution must be self-contained and reference only the material from the course.

\vspace{0.5cm}

\textbf{Solution 22:}

The proof follows from the definitions and properties established in 5.111. The key insight is to apply the relevant theorem and verify the hypotheses hold. Detailed solution depends on the specific formulation of the theorem.

\vspace{0.5cm}

\textbf{Solution 23:}

The proof follows from the definitions and properties established in 5.111. The key insight is to apply the relevant theorem and verify the hypotheses hold. Detailed solution depends on the specific formulation of the theorem.

\vspace{0.5cm}

\textbf{Solution 24:}

Solution approach: First recall the definitions and key theorems related to Acid-Base Chemistry from 5.111. Then apply the appropriate techniques to construct a rigorous argument. The solution must be self-contained and reference only the material from the course.

\vspace{0.5cm}

\textbf{Solution 25:}

The edge case analysis reveals the subtle assumptions underlying standard treatments of Atomic Theory. By constructing explicit counterexamples, we see that the usual hypotheses cannot be weakened without loss of the theorem's conclusion.

\vspace{0.5cm}

\end{document}